\section*{Capítulo \ref{ch:g10:storms-and-cyclones} --- Tormentas y ciclones tropicales}

\begin{solution}{exo:g10:storms-and-cyclones:1}
Una \omterm{def:g10:storms-and-cyclones:thunderstorm}{tormenta eléctrica} es una tormenta producida por un cumulonimbo, con
\omterm{def:g10:storms-and-cyclones:lightning}{rayos} y truenos, normalmente con chubascos intensos y rachas, y a veces con
granizo. Un \omterm{def:g10:storms-and-cyclones:lightning}{rayo} es una chispa eléctrica gigante que descarga regiones
cargadas de la tormenta. El trueno es el sonido del aire del canal del \omterm{def:g10:storms-and-cyclones:lightning}{rayo},
calentado de forma tan brusca que se dilata con una onda de choque.
\end{solution}

\begin{solution}{exo:g10:storms-and-cyclones:2}
Un \omterm{def:g10:storms-and-cyclones:tropical-cyclone}{ciclón tropical} es una tempestad de cientos de kilómetros de diámetro,
formada sobre un mar tropical cálido alrededor de un centro de bajas
presiones, con vientos que convergen en espiral y bandas de tormentas. El \omterm{def:g10:storms-and-cyclones:tropical-cyclone}{ojo}
es su centro, en calma y casi sin nubes; la \omterm{def:g10:storms-and-cyclones:tropical-cyclone}{pared del ojo} es el anillo de
tormentas altísimas que rodea el \omterm{def:g10:storms-and-cyclones:tropical-cyclone}{ojo}. Los vientos son más fuertes en la
\omterm{def:g10:storms-and-cyclones:tropical-cyclone}{pared del ojo}.
\end{solution}

\begin{solution}{exo:g10:storms-and-cyclones:3}
\num{70}~nudos: categoría 1. \num{120}~nudos: categoría 4.
\num{140}~nudos: categoría 5. Los dos últimos son huracanes mayores
(categoría 3 o superior).
\end{solution}

\begin{solution}{exo:g10:storms-and-cyclones:4}
$343 \times 9 \approx \qty{3.1}{km}$, y después $343 \times 15 \approx
\qty{5.1}{km}$. El segundo relámpago está más lejos: la tormenta se aleja.
\end{solution}

\begin{solution}{exo:g10:storms-and-cyclones:5}
El \omterm{def:g10:storms-and-cyclones:lightning}{rayo} tiende a caer en el punto más alto de los alrededores: un árbol
aislado en medio de un campo es ese punto, y la corriente de la descarga
baja por él y se extiende por el suelo a su alrededor. Hay que entrar en un
edificio o en un coche cerrado y quedarse dentro hasta bastante después del
último trueno.
\end{solution}

\begin{solution}{exo:g10:storms-and-cyclones:6}
Un nudo son $1852/3600 \approx \qty{0.514}{m/s}$. \num{64}~nudos: $64 \times
0.514 \approx \qty{32.9}{m/s}$, y $32.9 \times 3.6 \approx
\qty{119}{km/h}$. \num{137}~nudos: $\qty{70.5}{m/s}$, unos
\qty{254}{km/h} (la tabla oficial, redondeada a través de otra unidad, da
\qty{252}{km/h}).
\end{solution}

\begin{solution}{exo:g10:storms-and-cyclones:7}
$44 \times 86\,400 \approx \num{3.8e6}$ \omterm{def:g10:storms-and-cyclones:lightning}{rayos} al día, y $3.8\times10^6
\times 365.25 \approx \num{1.4e9}$, unos 1400 millones al año.
\end{solution}

\begin{solution}{exo:g10:storms-and-cyclones:8}
$\Delta p = 1010 - 935 = \qty{75}{hPa}$: un abombamiento de unos \qty{75}{cm}.
\end{solution}

\begin{solution}{exo:g10:storms-and-cyclones:9}
Se forman sobre los océanos tropicales cálidos, sobre todo entre
\qty{10}{\degree} y \qty{30}{\degree} de latitud en ambos hemisferios, con
muy pocos en el Atlántico Sur. Derivan hacia el oeste y luego se curvan hacia
los polos, donde se debilitan sobre aguas más frías. No aparecen sobre el
ecuador porque allí la rotación de la Tierra no imprime ningún giro al aire
que entra: sin él, el aire fluye directamente hacia las bajas presiones y
las rellena, y no puede nacer ninguna tempestad giratoria.
\end{solution}

\begin{solution}{exo:g10:storms-and-cyclones:10}
$1031 - 1.56 \times 70 \approx \qty{922}{hPa}$; $1031 - 1.56 \times 40
\approx \qty{969}{hPa}$.
\end{solution}

\begin{solution}{exo:g10:storms-and-cyclones:11}
Sobre tierra firme, la tempestad pierde su aporte de agua que se evapora, y
con él el \omterm{def:g10:earth-energy-budget:latent-heat}{calor latente} que la mueve; además, la superficie rugosa frena los
vientos. Sobre agua fría, la evaporación es mucho más débil, porque el aire
frío contiene mucho menos vapor de agua en la saturación: el aporte de \omterm{def:g10:earth-energy-budget:latent-heat}{calor
latente} disminuye.
\end{solution}

\begin{solution}{exo:g10:storms-and-cyclones:12}
A \qty{0}{\degreeCelsius}: $\sqrt{1.4 \times 8.314 \times 273.15/0.02897}
\approx \qty{331}{m/s}$; a \qty{30}{\degreeCelsius}, con $T =
\qty{303.15}{K}$: \qty{349}{m/s}. Para \qty{9}{s}: $9 \times 331 \approx
\qty{2980}{m}$ y $9 \times 349 \approx \qty{3140}{m}$, una diferencia de unos
\qty{160}{m}, alrededor de un \qty{5}{\percent}. La regla da \qty{3}{km}:
basta para saber lo cerca que está la tormenta, que es para lo único que
sirve.
\end{solution}

\begin{solution}{exo:g10:storms-and-cyclones:13}
$\Delta p = 1010 - 940 = \qty{70}{hPa}$: el abombamiento es de unos
\qty{0.7}{m}, y la parte del viento, $4.0 - 0.7 = \qty{3.3}{m}$, más de
cuatro veces más. Con la marea alta, el agua llega a $4.0 + 1.5 = \qty{5.5}{m}$ por encima del nivel medio del mar.
\end{solution}

\begin{solution}{exo:g10:storms-and-cyclones:14}
\omterm{def:g10:storms-and-cyclones:tornado}{Tornado}: una vuelta mide $2\pi \times 100 \approx \qty{628}{m}$, que se
recorren en $628/90 \approx \qty{7}{s}$. Huracán: $2\pi \times 30\,000 \approx
\qty{188}{km}$, que se recorren en $188\,000/60 \approx \qty{3140}{s}$, unos \qty{52}{min}. El \omterm{def:g10:storms-and-cyclones:tornado}{tornado} gira unas 450 veces más
deprisa: es un torbellino pequeño y muy rápido; el huracán, uno enorme y más
lento.
\end{solution}

\begin{solution}{exo:g10:storms-and-cyclones:15}
$24/3888 \approx \qty{0.6}{\percent}$; $414/3888 \approx
\qty{10.6}{\percent}$; $121/3888 \approx \qty{3.1}{\percent}$. Si la latitud
no importara, alrededor del \qty{9}{\percent} de las tempestades, unas 350,
nacerían a menos de \qty{5}{\degree} del ecuador; solo lo hacen 24, unas
quince veces menos. Cerca del ecuador, la rotación de la Tierra es demasiado
débil para poner en marcha el giro.
\end{solution}

\begin{solution}{pb:g10:storms-and-cyclones:1}
\textbf{1.} $\pi \times (665\,000)^2 \approx \qty{1.39e12}{m^2}$.

\textbf{2.} $0.015 \times 1.39\times10^{12} \approx \qty{2.08e10}{m^3}$.

\textbf{3.} Un kilómetro cúbico son \qty{e9}{m^3}: unos \qty{21}{km^3}.

\textbf{4.} $2.08\times10^{10} \times 1000 \approx \qty{2.08e13}{kg}$.

\textbf{5.} $2.08\times10^{13} \times 2.5\times10^6 \approx \qty{5.2e19}{J}$.

\textbf{6.} $5.2\times10^{19}/86\,400 \approx \qty{6.0e14}{W}$.

\textbf{7.} $6.0\times10^{14}/1.22\times10^{17} \approx 0.005$, alrededor del
\qty{0.5}{\percent}.

\textbf{8.} $240 \times 1.39\times10^{12} \approx \qty{3.3e14}{W}$. La
tempestad libera casi el doble de calor que la luz solar absorbida por toda
su superficie.

\textbf{9.} La tempestad no funciona con la luz solar que cae sobre ella:
libera calor almacenado en forma de vapor de agua, evaporado en una
superficie más amplia de mar cálido y arrastrado hacia dentro por los
vientos. La tempestad concentra energía reunida sobre una región mucho
mayor.

\textbf{10.} La misma masa que la lluvia, unos \qty{2.1e13}{kg} al día (el
vapor que entra tiene que renovarse tan deprisa como cae en forma de lluvia).

\textbf{11.} $6.0\times10^{14}/82 \approx \qty{7.3e12}{m^2}$, unas cinco
veces la superficie de la tempestad: la tempestad tiene que tomar vapor de
una zona mucho más amplia, y sus vientos hacen que el mar que hay debajo se
evapore mucho más deprisa que la media.

\textbf{12.} \qty{24.1}{g/kg} a \qty{28}{\degreeCelsius} frente a
\qty{14.7}{g/kg} a \qty{20}{\degreeCelsius}: $24.1/14.7 \approx 1.64$, un
\qty{64}{\percent} más.

\textbf{13.} El agua cálida se evapora deprisa y el aire cálido contiene
mucho más vapor: el aire que entra lleva mucho más \omterm{def:g10:earth-energy-budget:latent-heat}{calor latente} que
liberar y, una vez saturado, sube libremente. Sobre mares más fríos, el
aporte es demasiado pequeño para mover el motor.

\textbf{14.} Sobre tierra firme no hay mar que evaporar: el aporte de \omterm{def:g10:earth-energy-budget:latent-heat}{calor
latente} se corta, y el suelo rugoso frena los vientos.

\textbf{15.} $1.5\times10^{12} \times 86\,400 \approx \qty{1.3e17}{J}$.

\textbf{16.} $96 \times 1852/3600 \approx \qty{49.4}{m/s}$.

\textbf{17.} $\frac12 \times 1.16 \times 49.4^2 \approx \qty{1.4e3}{J}$.

\textbf{18.} $5.2\times10^{19}/1.3\times10^{17} = 400$: el \omterm{def:g10:earth-energy-budget:latent-heat}{calor latente} es
unas 400 veces la energía que reciben los vientos.

\textbf{19.} La mayor parte calienta el aire de las capas altas de la
tempestad, que se extiende por arriba y cede el calor al espacio en forma de
radiación infrarroja; los vientos pierden su propia energía por rozamiento
con el mar y con el suelo.

\textbf{20.} $1/400 = \qty{0.25}{\percent}$ (el programa da un cociente de
402): \textbf{alrededor de una parte de cada 400 del calor liberado por la
condensación se convierte en viento}.
\end{solution}
